\documentclass[11pt]{article}  % Reduce font size to 10pt
\usepackage[a4paper, total={6in, 8in}]{geometry}  % Reduce margins
\usepackage{amsthm}
\usepackage{amsfonts}
\usepackage{graphicx}
\usepackage{amssymb}
\usepackage{authblk}  % To manage author affiliations correctly
\usepackage[usenames]{xcolor}
\usepackage{lastpage}
\usepackage[utf8]{inputenc}
\usepackage{indentfirst}
\usepackage{amsfonts,amscd}
\usepackage{empheq}
\usepackage{latexsym}
\usepackage{graphicx,graphics,epsfig}
\usepackage{hyperref}

\usepackage{dsfont}

\urlstyle{same}

\renewcommand{\baselinestretch}{1.0}  % Tighten line spacing

\newtheorem{theorem}{Theorem}
\newtheorem{corollary}[theorem]{Corollary}
\newtheorem{lemma}[theorem]{Lemma}
\newtheorem{proposition}[theorem]{Proposition}
\newtheorem{conjecture}[theorem]{Conjecture}

\theoremstyle{definition}
\newtheorem{definition}[theorem]{Definition}
\newtheorem{remark}[theorem]{Remark}
\newtheorem{example}[theorem]{Example}

\makeatletter
\global\let\@jointwith\relax
\global\let\@affiliation\relax
% Redefine the \jointwith and \affiliation commands
\def\jointwith#1{\gdef\@jointwith{(Joint work with #1)}}
\def\affiliation#1{\gdef\@affiliation{#1}}

\newcommand{\Keywords}[1]{\par\noindent{\bfseries Keywords}: #1}
\newcommand{\MSC}[1]{\par\noindent{\bfseries MSC}: #1}

\renewcommand\maketitle{\par
  \begingroup
    \renewcommand\thefootnote{\@fnsymbol\c@footnote}%
    \def\@makefnmark{\rlap{\@textsuperscript{\normalfont\@thefnmark}}}%
    \long\def\@makefntext##1{\parindent 1em\noindent
            \hb@xt@1.8em{%
                \hss\@textsuperscript{\normalfont\@thefnmark}}##1}%
    \if@twocolumn
      \ifnum \col@number=\@ne
        \@maketitle
      \else
        \twocolumn[\@maketitle]%
      \fi
    \else
      \newpage
      \global\@topnum\z@
      \@maketitle
    \fi
    \thispagestyle{plain}\@thanks
  \endgroup
  \setcounter{footnote}{0}%
  \global\let\thanks\relax
  \global\let\@jointwith\relax
  \global\let\@affiliation\relax
  \global\let\maketitle\relax
  \global\let\@maketitle\relax
  \global\let\@thanks\@empty
  \global\let\@author\@empty
  \global\let\@date\@empty
  \global\let\@title\@empty
  \global\let\title\relax
  \global\let\author\relax
  \global\let\date\relax
  \global\let\and\relax
}
\def\@maketitle{%
  \newpage
  \null
  \vskip 2em%
  \begin{center}%
    {\LARGE \@title \par}%
    \vskip 1.5em%
    {\large
      \lineskip .5em%
      \begin{tabular}[t]{c}%
        \@author \\
        {\sc\small \@affiliation} \\
        {\small \@jointwith}
      \end{tabular}\par}%
    \vskip 1em%
  \end{center}%
  \par
  \vskip 1.5em}
\makeatother
\renewcommand*{\thefootnote}{\fnsymbol{footnote}}

\begin{document}

\title{Bloch-Wigner sequence and linear groups over finite rings}
\author{Abraham Rojas Vega}
\affiliation{Universidade de São Paulo}
\jointwith{Behrooz Mirzaii}
\maketitle

The Bloch-Wigner exact sequence is an important result in K-theory that connects the third homology group of $\operatorname{SL}_2(\mathbb{C})$ with the Bloch group, which encodes information about scissor congruences of polyhedra in hyperbolic geometry. These concepts provide deep insights into the interplay between algebra, geometry, and topology. 

Hutchinson, Mirzaii, and Pérez have extended the Bloch-Wigner sequence to other classes of commutative rings, like local rings \cite{hut2013,tesiselvis}.
In this work, we apply these results to study the low-dimensional homology groups of $\operatorname{SL}_2$ over finite commutative rings. We also discuss other achievements in the area that have been obtained using these methods.

\vline

\Keywords{Bloch group, Finite rings,}

\begin{thebibliography}{99}
\bibitem{hut2013}
\textsc{Hutchinson, K.}, A Bloch-Wigner sequence for SL2, \textit{Journal of K-Theory} \textbf{12}(2013), No.~1, 15--68.
\bibitem{tesiselvis}
\textsc{Torres Pérez, E.}, Scissors Congruence Group and the Third Homology of SL2, \textit{Tese de doutorado, Universidade de São Paulo} (2024)
\end{thebibliography}

\end{document}